% Modernised reading — fonds Alexandre Grothendieck, shelfmark 156-2, pages 1-18.
%
% This is an INTERPRETATION, not a transcription. Everything the critical
% apparatus carried moves into footnotes. In any dispute of substance the
% transcription governs: whoever cites Grothendieck must cite batch-01.fr.tex.
%
% Pass: Opus 5 (claude-opus-5), 2026-09-13 — first pass, unchecked against the
% pages by a human. The folder was read whole; it holds one batch, pages 1-18,
% all of them mathematics, dated 6 June 1986 (page 1) and 7 June (page 14).
%
% What the folder is. Chapter II (« GF II ») of Vers une géométrie des formes:
% the first steps of a topology of « failles » — a space with two ends,
% cut by « lieux » — and of « mailles » — a space cut open along a two-sided
% « nœud » into a faille. The « réseaux » of the chapter title are never
% reached: the folder builds the two building blocks and stops at the
% compatibility of two mailles.
%
% Notation fixed for the whole document. X/A the space X with A contracted to
% a point a. V ≅ C(S) a conical neighbourhood, S its base, V_t = tV, W_t the
% preimage of V_t in X, S_t its boundary. F = (X, {A,B}) a faille, U =
% X∖(A∪B), X' the reduced faille (A, B contracted to a, b). L(F) the set of
% lieux; for an interior lieu C, X∖C = U_{A,C} ⊔ U_{B,C}, X_{A,C} =
% closure(U_{A,C}) ∪ C, F_{A,C} = (X_{A,C}, {A,C}). Comp(Z) the boolean ring
% of clopen subsets. Riv(A,X) the rives; Ã the cut of A; β = {ρ,ρ'} a
% bi-rive; X̄ = Dec_β(A,X) the space cut along (A,β). « Disjoints » for lieux
% is always the page's relation (defined page 8), never mere empty
% intersection; where the two could be confused the reading says the lieux
% are « emboîtés ».
%
% Substantive departures from the pages, each footnoted where it occurs:
% — page 2, NB (3): a closed submanifold is regularly immersed, via the normal
%   disc bundle; the reading says which cone and in what topology.
% — page 3: c') is stated with the closure of U_i (his margin), and the
%   equivalence c) ⇔ c') is proved one way in general and the other when U
%   has finitely many components, as he says.
% — page 5: the last contradiction is with the connectedness of X∖{a}; the
%   leaf writes X∖{c}.
% — page 6: the lemma needs Z = Z' ∪ Z''; supplied.
% — page 8: « disjoints » (his) implies C ∩ D = ∅ but not conversely; said.
% — page 9: the parenthesis read « totalement ordonnée » is taken to mean that
%   pairwise disjoint lieux form a chain, which is true; the set of all lieux
%   between two lieux is not totally ordered, and the reading says so.
% — page 10: the existence criterion and the uniqueness on the segment are
%   true of REDUCED failles; « ∞ de structures réduites » fails for finite
%   trees (k(k-1)/2); uniqueness for the segment is R. L. Moore's non-cut-point
%   theorem. The two retracted claims are recorded as retracted.
% — page 12: « homéomorphe à B^n » is read as: X ≅ B^n, ends and lieux
%   ≅ B^{n-1}; the sketch's D^1 is B^1.
% — page 15: ρ + ρ' = 1 forces ρ ≠ ρ' as soon as A is not open (the page says
%   A ≠ ∅); (*) is injective always, bijective for X paracompact (Godement).
% — page 16: X̃ is used without definition; the reading describes it and says
%   that the precise definition is on neither the page nor here.
%
% Modern names given and footnoted as ours: voisinage conique / voisinage
% régulier (J. H. C. Whitehead 1939), espaces localement coniques
% (Siebenmann 1972), cobordisme et sa composition (Thom 1954; the categorical
% form after 1986), théorème de R. L. Moore sur les points de non-coupure
% (1920), plongement bicollier (M. Brown 1962) and théorème de Schoenflies
% généralisé (Brown 1960, Mazur 1959), dualité de Stone (1936) and spectre de
% Stone relatif, hypersurface bilatère non séparante.
%
% What the folder announces and never establishes: the equivalences of
% categories of pages 7-9 (failles with lieux ↔ chains of failles glued by
% homeomorphisms); the claim of page 7 that every lieu of a component faille
% is a lieu of F; the symmetry of the compatibility of two mailles (page 18),
% on which the text stops; the « réseaux » of the title.
\documentclass[11pt,a4paper]{article}
% Le préambule commun à toutes les transcriptions du fonds Grothendieck.
%
% Il tient en une idée : **l'incertitude se compose, elle ne se résout pas.**
% Une transcription qui lisse un mot douteux a détruit la seule chose pour
% laquelle on la faisait — un lecteur doit pouvoir savoir, à chaque endroit,
% ce qui était sur le feuillet et ce qui a été ajouté. Les six macros qui
% suivent sont donc l'appareil critique, et non de la décoration.
%
% Les mêmes commandes sont reconnues par `scripts/render.mjs`, qui produit la
% vue de lecture du site. Une macro ajoutée ici doit l'être aussi là-bas,
% faute de quoi le rendu échouera bruyamment — ce qui est voulu : un rendu
% silencieusement faux est pire qu'un rendu absent.

\ProvidesPackage{grothendieck}[2026/08/08 Transcriptions du fonds Grothendieck]

\usepackage{amsmath,amssymb,amsthm}
\usepackage{mathrsfs}
\usepackage{stmaryrd}
% Les diagrammes commutatifs sont partout dans ce fonds : c'est la langue
% même de la géométrie algébrique de Grothendieck, et une transcription qui
% les rendrait en prose ne transcrirait rien.
\usepackage{tikz-cd}
% Français par défaut — c'est la langue des feuillets — mais l'anglais est
% chargé aussi : la traduction et le résumé sont des documents anglais, et la
% typographie française (espaces avant les deux-points) y serait une faute.
% Ces documents basculent par \selectlanguage{english} en tête de corps.
\usepackage[english,french]{babel}
\usepackage{fontspec}
\usepackage{geometry}
\geometry{a4paper,margin=25mm,marginparwidth=28mm}
\usepackage{marginnote}
\usepackage{xcolor}
% ulem, not soul. `soul`'s \st and \ul are built on a character-by-character
% scan that XeLaTeX's font machinery breaks: the document fails with an
% undefined control sequence rather than degrading. ulem does the same two jobs
% and is Unicode-engine safe. `normalem` stops it hijacking \emph, which the
% transcriptions use for Grothendieck's own emphasis.
\usepackage[normalem]{ulem}
\usepackage{hyperref}
\hypersetup{colorlinks=true,linkcolor=black,urlcolor=black,pdfborder={0 0 0}}

% --- Métadonnées du lot -----------------------------------------------------
% Reprises telles quelles par le script de rendu, qui les affiche en tête de
% la vue de lecture. Elles fixent ce que le fichier prétend couvrir : sans
% elles, une transcription flotte sans point d'attache dans un fonds de seize
% mille feuillets.

\newcommand{\folder}[1]{\gdef\grothFolder{#1}}
\newcommand{\batch}[1]{\gdef\grothBatch{#1}}
\newcommand{\leaves}[2]{\gdef\grothFirst{#1}\gdef\grothLast{#2}}
\newcommand{\foldertitle}[1]{\gdef\grothFoldertitle{#1}}
\gdef\grothFolder{?}\gdef\grothBatch{?}\gdef\grothFirst{?}\gdef\grothLast{?}\gdef\grothFoldertitle{}

% --- Le résumé --------------------------------------------------------------
%
% Il ouvre la lecture modernisée, sous le titre « Résumé », et il est composé
% en petit. Ce n'est pas un ornement : c'est le seul passage du document qui
% ne soit pas des mathématiques, et un lecteur doit pouvoir le reconnaître
% avant de l'avoir lu — puis le sauter s'il connaît déjà le sujet, ou s'y
% arrêter s'il ne le connaît pas. Le filet qui le suit marque la reprise du
% corps.

\newenvironment{resume}
  {\small\setlength{\parskip}{0.45em}}
  {\par\medskip\hrule\medskip}

% --- Mots-clés ---------------------------------------------------------------
%
% En anglais, en clôture du résumé de la lecture modernisée : le vocabulaire
% moderne sous lequel chercher ce que les feuillets construisent. Le manifeste
% du site (`npm run manifest`) les extrait pour étiqueter la cote entière —
% cette ligne est la source unique des tags, il n'y a pas de fichier à côté.
\newcommand{\keywords}[1]{\par\smallskip\noindent
  {\footnotesize\textsc{Keywords} --- \textit{#1}}\par}

% --- L'appareil critique ----------------------------------------------------

% Le numéro de feuillet, en marge. C'est l'ancre qui permet de régler un
% désaccord sur une formule en regardant *un* feuillet plutôt qu'une chemise
% de six cents. Le site s'en sert aussi pour tourner les pages du fac-similé
% quand on fait défiler la transcription.
\newcommand{\leaf}[1]{%
  \marginnote{\footnotesize\color{gray!70}\textsf{#1}}%
  \label{leaf:#1}\ignorespaces}

% Illisible : on ne devine pas. Un mot inventé qui se lit comme les autres est
% le pire résultat possible.
\newcommand{\ill}{\textcolor{red!60!black}{[\,\ldots\,]}}

% Lecture douteuse : le mot est donné, et signalé comme incertain.
\newcommand{\uncertain}[1]{\uline{#1}}

% Ajout de l'éditeur — une lettre manquante, un mot sous-entendu.
\newcommand{\add}[1]{\textcolor{blue!50!black}{[#1]}}

% Biffé par Grothendieck lui-même. Ce qu'il a rayé fait partie du document :
% une rature dit souvent où la pensée a changé de direction.
\newcommand{\struck}[1]{\sout{#1}}

% Note du transcripteur, sur l'état matériel du feuillet.
\newcommand{\note}[1]{\par\smallskip{\small\color{gray!60!black}\textit{[#1]}}\par\smallskip}

% Note marginale de Grothendieck, distincte du corps du texte.
\newcommand{\marginal}[1]{\par\smallskip\noindent\hspace{1em}%
  {\small\color{gray!60!black}#1}\par\smallskip}

% --- Titre ------------------------------------------------------------------

\AtBeginDocument{%
  \begin{center}
    {\large\bfseries Fonds Alexandre Grothendieck --- cote n\textsuperscript{o}~\grothFolder}\\[2pt]
    {\itshape\grothFoldertitle}\\[4pt]
    {\small feuillets \grothFirst--\grothLast\ (lot \grothBatch)}\\[2pt]
    {\footnotesize\color{gray!60!black}Transcription établie d'après le fac-similé numérisé
     par l'Université de Montpellier.\\
     Les lectures incertaines sont soulignées, les illisibles notées [\,\ldots\,],
     les ajouts entre crochets.}
  \end{center}
  \vspace{1em}}

\endinput


\folder{156-2}
\pages{1}{18}
\dating{1986}
\watermark{Édition de démonstration\\interprétation personnelle de l'œuvre}
\foldertitle{[Chapitre] II. Réalisations topologiques des réseaux : notes manuscrites (06/06/1986) — lecture modernisée du dossier entier}

\begin{document}

\section*{Résumé}

\begin{resume}
Un segment a deux bouts, et tout point intérieur le coupe en deux segments
plus petits. Un cercle n'a pas de bouts, mais si on le coupe en un point, on
obtient un segment dont les deux bouts sont deux copies du point de coupure.
Ces deux gestes --- couper un objet entre ses bouts, et ouvrir une boucle en
la coupant --- sont ce que ces dix-huit feuillets, écrits les 6 et 7 juin 1986,
cherchent à formuler pour des espaces topologiques quelconques, sans
supposer qu'ils soient des variétés ni même triangulés.

La première notion est la \emph{faille} : un espace $X$ muni de deux parties
fermées disjointes $A$ et $B$, ses bouts, qui se comporte « comme un segment »
du point de vue de la connexité --- ôter un bout ne déconnecte pas l'espace,
et près de chaque bout l'espace a la forme d'un cône. Un cylindre
$S^{1} \times [0,1]$ avec ses deux cercles de bord en est l'exemple type, une
boule avec deux points de son bord en est un autre. Dans une faille, un
\emph{lieu} est une partie qui la coupe en deux failles, l'une du côté de $A$,
l'autre du côté de $B$ : un cercle horizontal dans le cylindre. Les lieux
s'ordonnent de $A$ vers $B$, et découper une faille par plusieurs lieux
revient à recoller bout à bout plusieurs failles. Qui connaît les cobordismes
reconnaîtra leur composition, mais ici rien n'est lisse et l'on ne parle que
d'espaces.

La seconde notion est la \emph{maille}. Près d'une partie fermée $A$ de $X$,
l'espace privé de $A$ peut avoir plusieurs « côtés » : ce sont les
\emph{rives} de $A$. Une hypersurface bilatère a deux côtés ; l'âme d'un ruban
de Möbius n'en a qu'un, bien qu'elle coupe le ruban en deux près de chacun de ses
points. Choisir une partition
des rives en deux, puis couper $X$ le long de $A$ en dédoublant $A$, donne un
nouvel espace qui a deux bouts, copies de $A$ ; si c'est une faille, on dit que
$A$ est le \emph{nœud} d'une maille. Le cercle coupé en un point est la plus
simple ; un anneau coupé le long d'un rayon en est une autre.

Le titre annonce des « réseaux », et l'on devine l'intention : bâtir des
espaces en recollant failles et mailles, comme on bâtit un graphe avec des
arêtes et des boucles. Les feuillets n'y arrivent pas ; ils s'arrêtent sur la
question de savoir quand deux mailles d'un même espace sont compatibles. Mais
les briques sont posées, et plusieurs remarques en marge --- « Faux ! idiot ! »,
« Faux … déjà dans le cas des segments ! » --- montrent un auteur qui
essaie des énoncés et les retire aussitôt qu'un exemple simple les défait.

\keywords{conelike neighbourhood, regular neighbourhood, cobordism, non-cut point, bicollared embedding, generalized Schoenflies theorem, Stone duality, two-sided hypersurface, cutting along a subspace}
\end{resume}

\section*{Le fil du dossier, et les conventions}

\begin{enumerate}
\item[1.] \emph{Immersion régulière} (pages 1 et 2).
\item[2.] \emph{Failles topologiques} (pages 2 à 4).
\item[3.] \emph{Lieux, et découpage d'une faille} (pages 4 à 7).
\item[4.] \emph{L'ordre des lieux} (pages 8 et 9).
\item[5.] \emph{Existence, et conditions sur les lieux} (pages 10 à 13).
\item[6.] \emph{Rives et bi-rives} (pages 14 et 15), sous le titre « 2) Boucle
topologique » daté du 7 juin.
\item[7.] \emph{Mailles} (pages 16 à 18).
\end{enumerate}

Les feuillets portent sa pagination 1 à 18, qui coïncide avec celle des
archivistes. La page 1 porte en angle « GF II » : c'est le chapitre II d'un
texte dont le dossier 156-1, « Vers une géométrie des formes (topologiques) »,
daté de la veille, est le chapitre I.

\emph{Conventions.} Pour une partie fermée $A$ d'un espace $X$, $X/A$ est
l'espace obtenu en contractant $A$ en un point $a$. $\mathrm{C}(S) = S \times
[0,1] / S \times \{0\}$ est le cône sur $S$, de sommet $o$. Une faille est notée
$F = (X, \{A,B\})$, $U = X \smallsetminus (A \cup B)$, et $X'$ la faille réduite
obtenue en contractant $A$ et $B$ en $a$ et $b$. Pour un lieu intérieur $C$,
$X \smallsetminus C = U_{A,C} \sqcup U_{B,C}$, $X_{A,C} = \overline{U_{A,C}} \cup
C$ et $F_{A,C} = (X_{A,C}, \{A,C\})$. $\mathrm{Comp}(Z)$ est l'anneau booléen des
parties ouvertes et fermées de $Z$. Le mot « disjoints », appliqué à deux lieux,
a toujours le sens fixé à la page 8, plus fort que l'intersection vide.

\emph{Ce que le dossier annonce et n'établit pas} : les équivalences de
catégories des pages 7 à 9 entre failles munies de lieux et chaînes de failles
recollées ; l'affirmation de la page 7 que tout lieu d'une faille composante est
un lieu de la faille ; la symétrie de la compatibilité de deux mailles, sur
laquelle le texte s'arrête (page 18) ; les réseaux du titre.

\pagerange{1}{2}
\section*{Immersion régulière (pages 1 et 2)}

Soit $X$ un espace topologique et $A \subset X$ une partie fermée non
vide\footnote{« non vide » est ajouté au-dessus de la ligne ; la remarque (2) de
la page 2 examine pourtant le cas $A = \emptyset$, voir plus bas.}. Si $X'
\subset X$ contient $A$, l'application $X'/A \to X/A$ identifie $X'/A$ à un
sous-espace de $X/A$ : les fermés de $X'/A$ correspondent aux fermés $Y'$ de $X'$
tels que $Y' \cap A \in \{\emptyset, A\}$, et ce sont les traces sur $X'$ des
fermés de $X$ ayant la même propriété.

On dit que $A$ est \emph{régulièrement immergé} dans $X$ si le point $a$ possède
dans $X/A$ un voisinage fermé $V$ muni d'un homéomorphisme $h : V \simeq
\mathrm{C}(S)$ envoyant $a$ sur le sommet, où $S$ est un espace\footnote{La page
qualifie $S$ d'un mot lu « compact » avec doute. Rien de ce qui suit ne s'en
sert, sinon pour que le cône $\mathrm{C}(S)$, pris avec la topologie quotient,
soit un voisinage raisonnable de son sommet ; on ne l'impose pas.}, et tel que
$V \smallsetminus B$ soit ouvert dans $X/A$, où $B = h^{-1}(S \times \{1\})$ est
le bord de $V$\footnote{La lettre $S$ est écrite en surcharge sur un $B$ : le
bord de $V$ et la base du cône se correspondent par $h$.}. La structure conique
fournit les $V_{t} = h^{-1}(\mathrm{C}_{t}(S))$, $0 < t \leq 1$, tous homéomorphes
à $V$, qui forment un système fondamental de voisinages de $a$. Leurs images
réciproques $W_{t}$ dans $X$ forment alors un système fondamental de voisinages
de $A$ --- tout ouvert de $X$ contenant $A$ est saturé pour la contraction, donc
provient d'un ouvert de $X/A$. Le bord $S_{t}$ de $W_{t}$ est fermé, et
\[
X \smallsetminus S_{t} = W_{t}^{\circ} \sqcup (X \smallsetminus W_{t})
\]
est une somme topologique de l'intérieur et de l'extérieur de $W_{t}$.

La notion est locale : elle ne dépend que d'un voisinage de $A$. Elle a encore un
sens pour $A = \emptyset$, où $X/\emptyset = X \sqcup \{a\}$ et le cône est sur
l'espace vide. Elle est vraie, dit la page, pour une sous-variété d'une
variété\footnote{L'énoncé de la page est en partie illisible. Pour une
sous-variété fermée lisse $A$ d'une variété lisse $X$, sans bord, le fibré en
disques normaux $D(\nu)$ est un voisinage fermé de $A$, cylindre de
l'application $S(\nu) \to A$ ; son quotient $D(\nu)/A$ est le cône
$\mathrm{C}(S(\nu))$ sur le fibré en sphères, avec la topologie quotient. C'est
l'hypothèse sous laquelle on la tient pour vraie ici.}.

Le nom actuel le plus proche est celui de \emph{voisinage conique} ou de
\emph{voisinage régulier}\footnote{Les voisinages réguliers de J.~H.~C.~Whitehead
(1939) en topologie PL, les espaces localement coniques de L.~Siebenmann (1972),
et plus généralement les espaces stratifiés, sont les cadres où cette condition
est d'usage. Grothendieck en connaissait l'esprit : l'\emph{Esquisse d'un
programme} (1984) réclame une « topologie modérée » précisément pour ce genre
d'objets. Les feuillets ne citent personne.}. On notera que la condition porte
sur $X/A$ et non sur $X$ : $A$ lui-même peut être un espace sauvage, seul compte
le voisinage une fois $A$ écrasé.

\pagerange{2}{4}
\section*{Failles topologiques (pages 2 à 4)}

\textbf{Définition.} Une \emph{faille topologique} est un couple $(X, \{A,B\})$
formé d'un espace $X$ et d'une paire de parties de $X$ tels que :
\begin{enumerate}
\item[a)] $A$ et $B$ sont fermés et non vides, et les points $a \in X/A$,
$b \in X/B$ ne sont pas isolés ;
\item[b)] $A \cap B = \emptyset$ ;
\item[c)] $(X \smallsetminus A)/B$ et $(X \smallsetminus B)/A$ sont connexes ;
\item[d)] $A$ et $B$ sont régulièrement immergés.
\end{enumerate}

$A$ et $B$ sont les \emph{bouts} de la faille. Comme $a$ n'est pas isolé, $A$
n'est pas ouvert, et de même pour $B$ ; donc $X \neq A \cup B$, et
$U = X \smallsetminus (A \cup B)$ n'est pas vide.

La condition c) admet une variante :
\begin{enumerate}
\item[c')] l'adhérence de toute composante connexe $U_{i}$ de $U$ rencontre $A$
et rencontre $B$.\footnote{Le corps de la page écrit « $U_{i}$ rencontre $B$ et
rencontre $A$ », ce qui est impossible puisque $U_{i} \subset U$ ; la marge
corrige en $\overline{U_{i}}$.}
\end{enumerate}
On a c') $\Rightarrow$ c), et la réciproque est vraie si $U$ n'a qu'un nombre
fini de composantes connexes, ce qui est le cas « modéré »\footnote{Les
démonstrations ne sont pas sur la page. c') $\Rightarrow$ c) : une partie
ouverte et fermée de $(X \smallsetminus A)/B$ qui ne contient pas $b$ contient
chaque $U_{i}$ qu'elle rencontre, donc son adhérence, donc un point de $B$ ---
contradiction. Réciproque : si $U$ n'a qu'un nombre fini de composantes, chacune
est ouverte et fermée dans $U$ ; une composante dont l'adhérence ne rencontre pas
$B$ serait ouverte et fermée dans $X \smallsetminus A$ et saturée, et
déconnecterait $(X \smallsetminus A)/B$.}.

La faille est \emph{réduite} si $A$ et $B$ sont des points $a$ et $b$. En
contractant les deux bouts on obtient $X' = X/A/B$, et $(X, \{A,B\})$ est une
faille si et seulement si $(X', \{a,b\})$ en est une. Pour une faille réduite,
les conditions se lisent : $a$ et $b$ ne sont pas isolés, $a \neq b$,
$X' \smallsetminus \{a\}$ et $X' \smallsetminus \{b\}$ sont connexes, et $a$, $b$
ont des voisinages coniques. Une faille réduite est donc connexe, puisque
$X' \smallsetminus \{b\}$ est connexe et dense. En revanche
$X' \smallsetminus \{a,b\}$ ne l'est pas nécessairement : le cercle muni de deux
points en est l'exemple\footnote{La marge écrit « mais pas néc. », d'une
abréviation lue avec doute ; la sphère $S^{2}$ munie de deux points, dont le
complémentaire est connexe, montre qu'on ne peut pas dire mieux.}.

Qu'on puisse contracter les bouts sans rien perdre est la raison pour laquelle
presque tout ce qui suit se démontre « on peut supposer $A$, $B$ réduits à des
points ».

\pagerange{4}{7}
\section*{Lieux, et découpage d'une faille (pages 4 à 7)}

\subsection*{Définition}

\begin{quote}
\textbf{Définition.} Un \emph{lieu} de la faille $(X, \{A,B\})$ est une partie
fermée $C$ de $X$ telle que ou bien $C \in \{A, B\}$, ou bien $C \subset U$ et :
il existe une décomposition de $X \smallsetminus C$ en somme topologique de deux
sous-espaces
\[
X \smallsetminus C = U_{A,C} \sqcup U_{B,C}, \qquad A \subset U_{A,C}, \quad
B \subset U_{B,C},
\]
avec $U_{A,C}/A$ et $U_{B,C}/B$ connexes.
\end{quote}

Autrement dit, dans la faille réduite, $X' \smallsetminus C$ a exactement deux
composantes connexes, l'une contenant $a$ et l'autre $b$ ; la décomposition est
donc unique\footnote{La page amorce « et ceci détermine $U_{A,C}$, $U_{B,C}$ de
façon unique, comme les composantes connexes de $X' \smallsetminus C$ qui
contiennent » et s'interrompt ; la marge le dit en toutes lettres. Une note
oblique, en partie raturée, ajoute quelque chose sur le cas où $A$ et $B$ sont
connexes.}. Un lieu intérieur est non vide, puisque $X'$ est connexe.

On pose $X_{A,C} = \overline{U_{A,C}} \cup C$ et $X_{B,C} = \overline{U_{B,C}}
\cup C$.

\subsection*{Les deux moitiés sont presque des failles}

\begin{quote}
\textbf{Proposition} (page 5). \emph{$(X_{A,C}, \{A,C\})$ satisfait les conditions
a), b), c) des failles, et de même $(X_{B,C}, \{B,C\})$.}
\end{quote}

On se ramène au cas où $A$, $B$, $C$ sont des points $a$, $b$, $c$. Comme
$U_{a,c}$ est ouvert et fermé dans $X \smallsetminus \{c\}$, son adhérence est
contenue dans $U_{a,c} \cup \{c\}$.

a) $a$ n'est pas isolé dans $X_{a,c}$, qui contient un voisinage de $a$ dans $X$.
Si $c$ l'était, $U_{a,c}$ serait fermé dans $X$ ; il est ouvert ; il serait
ouvert et fermé, non vide et distinct de $X$, contre la connexité de $X$.

c) $X_{a,c} \smallsetminus \{c\} = U_{a,c}$ est connexe par hypothèse. Si
$X_{a,c} \smallsetminus \{a\}$ ne l'était pas, il contiendrait une partie $P$
ouverte et fermée, non vide, ne contenant pas $c$. Fermée dans
$X_{a,c} \smallsetminus \{a\}$, qui est fermé dans $X \smallsetminus \{a\}$, et
ouverte dans l'ouvert $U_{a,c} \smallsetminus \{a\}$ de $X$, elle serait ouverte
et fermée dans $X \smallsetminus \{a\}$, qui ne serait pas connexe ---
absurde\footnote{La page conclut « $X \smallsetminus \{c\}$ ne serait pas
connexe », ce qui n'est pas une contradiction : $X \smallsetminus \{c\}$ est
déconnecté par définition d'un lieu. La contradiction est avec la connexité de
$X \smallsetminus \{a\}$, où l'argument vient de placer $P$.}.

b) est clair. Il ne manque, pour avoir une faille, que la condition d)
d'immersion régulière du nouveau bout $C$.

\subsection*{La condition qui manque}

\begin{quote}
\textbf{Lemme} (page 6). \emph{Soit $Z$ un espace réunion de deux parties fermées
$Z'$ et $Z''$ telles que $Z' \cap Z'' = \{c\}$. Pour que $c$ ait un voisinage
conique dans $Z$, il faut et il suffit qu'il en ait un dans $Z'$ et dans $Z''$.}
\end{quote}

L'hypothèse $Z = Z' \cup Z''$ n'est pas écrite, mais c'est ainsi que le lemme
sert\footnote{Sans elle l'énoncé n'a pas de sens. La démonstration n'est pas sur
la page. Suffisance : $\mathrm{C}(S') \vee_{o} \mathrm{C}(S'') =
\mathrm{C}(S' \sqcup S'')$. Nécessité : si $V \simeq \mathrm{C}(S)$ est un
voisinage de $c$ dans $Z$, la trace de $Z'$ sur $V \smallsetminus \{c\} \simeq S
\times {]0,1]}$ y est ouverte et fermée, donc de la forme $T \times {]0,1]}$ avec
$T$ ouvert et fermé dans $S$, et $V \cap Z' \simeq \mathrm{C}(T)$.}. Appliqué à
$X' / C = X_{A,C}/C \cup X_{B,C}/C$, il dit que si l'on ajoute à la définition
d'un lieu intérieur la condition « $C$ est régulièrement immergé dans $X$ », les
deux moitiés $F_{A,C}$ et $F_{B,C}$ sont des failles. C'est ce que la page
décide : on exige désormais cette condition.

\emph{Exemple.} Pour $t$ assez petit, les bords $S_{A,t}$ des voisinages coniques
$W_{t}$ de $A$ sont des lieux de la faille\footnote{Pour $t$ petit, $W_{t}$ évite
$B$, puisque les $W_{t}$ forment un système fondamental de voisinages de $A$ ;
$W_{t}^{\circ}/A$ est un cône ouvert, donc connexe, et
$(X \smallsetminus W_{t})/B$ est connexe parce qu'on peut y rendre le collier
$S \times {]t, 1[}$ à $(X \smallsetminus A)/B$ sans changer ses parties ouvertes
et fermées. La page renvoie à une page dont le numéro est illisible.}.

\subsection*{Découper et recoller}

Un lieu intérieur $C$ de $F = (X, \{A,B\})$ décompose la faille en deux failles
$F_{A,C}$ et $F_{B,C}$, et $X = X_{A,C} \cup_{C} X_{B,C}$. La page affirme que
tout lieu d'une faille composante est un lieu de $F$, et que la catégorie des
failles munies d'un lieu intérieur, avec les homéomorphismes, est équivalente à
celle des paires de failles munies d'un homéomorphisme entre un bout de l'une et
un bout de l'autre\footnote{Ces deux énoncés sont écrits sans démonstration, le
second dans une énumération biffée puis reprise en interligne. On ne les
démontre pas ici.}. C'est, sans aucune structure lisse, la forme de la
composition des cobordismes\footnote{Un cobordisme $W$ entre deux variétés
$\partial_{0} W$ et $\partial_{1} W$ est une faille $(W, \{\partial_{0}W,
\partial_{1}W\})$ dès que ses bords ont des colliers et que les conditions de
connexité sont satisfaites, et recoller deux cobordismes le long d'un bord commun
est l'opération ci-dessus. La notion de cobordisme remonte à R.~Thom (1954) ; sa
présentation comme composition dans une catégorie est postérieure à juin 1986
(M.~Atiyah, 1988). Le rapprochement est le nôtre.}.

\begin{quote}
\textbf{Lemme} (page 7). \emph{Soient $C \neq D$ deux lieux de $F$ distincts des
bouts. Alors $D$ est un lieu de $F_{B,C}$ si et seulement si $C$ est un lieu de
$F_{A,D}$ ; et, en échangeant les rôles de $A$ et $B$, $D$ est un lieu de
$F_{A,C}$ si et seulement si $C$ est un lieu de $F_{B,D}$.}
\end{quote}

Il suffit de prouver une implication, l'autre s'obtenant en échangeant $A$ et $B$,
$C$ et $D$ ; et l'on peut supposer $A$, $B$, $C$, $D$ réduits à des points. Si $c$
est un lieu de $X$ et $d$ un lieu de $X_{c,b}$, on a
\[
X = X_{a,c} \amalg_{c} X_{c,b}, \qquad X_{c,b} = X_{c,d} \amalg_{d} X_{d,b},
\qquad \text{donc} \qquad
X = \underbrace{X_{a,c} \amalg_{c} X_{c,d}}_{X_{a,d}} \amalg_{d} X_{d,b},
\]
et $c$ est un lieu de $X_{a,d}$\footnote{La page note $X_{c,b} = X_{b,c}$ par
définition. La démonstration s'arrête sur la décomposition, qui est bien le
fond de l'affaire : $X_{a,d} \smallsetminus \{c\}$ est la somme de
$X_{a,c} \smallsetminus \{c\}$ et de $X_{c,d} \smallsetminus \{c\}$, chacune
connexe une fois son bout contracté. Le lemme est réécrit à la page 12, puis
barré.}.

\pagerange{8}{9}
\section*{L'ordre des lieux (pages 8 et 9)}

\begin{quote}
\textbf{Définition} (page 8). Deux lieux $C$ et $D$ sont \emph{disjoints} si :
lorsque l'un d'eux est un bout, $C \neq D$ ; sinon, $C \neq D$ et $D$ est un lieu
de $F_{A,C}$ ou de $F_{B,C}$ --- ce qui, par le lemme de la page 7, équivaut à ce
que $C$ soit un lieu de $F_{B,D}$ ou de $F_{A,D}$. La relation est symétrique.
\end{quote}

Deux lieux disjoints en ce sens ont une intersection vide, mais la réciproque est
fausse\footnote{La page 6 annonce « disjoints au sens $C \cap D = \emptyset$ » pour
le seul cas où l'un des lieux est un bout, où les deux notions coïncident. Pour
deux lieux intérieurs, $D \cap C = \emptyset$ ne suffit pas : un lieu n'est pas
supposé connexe, et $D$ peut rencontrer à la fois $U_{A,C}$ et $U_{B,C}$.} : la
relation de la page 8 dit que les deux lieux sont \emph{emboîtés}, l'un du côté de
$A$ par rapport à l'autre.

Une faille \emph{orientée} est une faille dont on a choisi un bout pour origine ;
on écrit $F = (X, (A,B))$, et $\mathrm{or}(F) = A$, $\mathrm{ex}(F) = B$. Pour
un lieu $C$ on pose $\mathcal{F}_{C}$ l'ensemble des lieux de $F_{A,C}$ si $C$ est
intérieur, $\mathcal{F}_{A} = \{A\}$\footnote{La valeur pour $C = B$ est
illisible ; la seule qui rende la suite cohérente est l'ensemble de tous les
lieux de $F$.}, et l'on écrit
\[
C \leq D \iff \mathcal{F}_{C} \subset \mathcal{F}_{D} \iff C \in
\mathcal{F}_{D}.
\]
C'est une relation d'ordre --- l'antisymétrie vient de ce que $F_{A,C} =
F_{A,D}$ entraîne $C = D$, comme le dit la marge --- et deux lieux sont disjoints
si et seulement si $C < D$ ou $D < C$. Changer l'orientation remplace l'ordre par
l'ordre opposé.

\textbf{Proposition} (page 9). \emph{L'ensemble ordonné $\mathcal{L}$ des lieux
d'une faille orientée satisfait :}
\begin{enumerate}
\item[a)] \emph{il a un plus petit élément $\alpha = A$ et un plus grand élément
$\beta = B$, distincts ;}
\item[b)] \emph{pour tout $x \neq \beta$, l'ensemble des $y > x$ est filtrant
décroissant ;}
\item[b')] \emph{pour tout $x \neq \alpha$, l'ensemble des $y < x$ est filtrant
croissant ;}
\item[c)] \emph{si $x < y$, il existe $z$ avec $x < z < y$ (divisibilité).}
\end{enumerate}

La page énonce ces propriétés sans démonstration ; elles résultent de l'exemple
de la page 6\footnote{Pour b), si $y, y' > x$, les bords $S_{x,t}$ des petits
voisinages coniques de $x$ évitent le fermé $y \cup y'$ et sont des lieux compris
entre $x$ et chacun d'eux ; pour c), on applique le même argument dans la faille
$F_{x,y}$. C'est ici que l'exigence d'immersion régulière des lieux, ajoutée à la
page 6, est indispensable.}. L'ordre n'est pas total : deux cercles ondulés qui se
croisent sur un cylindre sont deux lieux non comparables. Ce qui est vrai, et
qu'on lit dans la parenthèse de la page 9, c'est qu'une famille de lieux deux à
deux disjoints est totalement ordonnée\footnote{La parenthèse est en grande
partie illisible ; seuls « forment une partie totalement ordonnée » se lisent. On
l'interprète dans le seul sens où l'énoncé est juste.}.

Enfin, découper une faille orientée par $n$ lieux intérieurs deux à deux disjoints
$x_{1} < \dots < x_{n}$ revient à se donner $n+1$ failles orientées recollées bout
à bout : la page affirme l'équivalence de cette catégorie avec celle des systèmes
\[
(F_{1}, \dots, F_{n+1};\; h_{21}, \dots, h_{n+1,n}), \qquad h_{i+1,i} :
\mathrm{ex}(F_{i}) \xrightarrow{\ \simeq\ } \mathrm{or}(F_{i+1}),
\]
de $n+1$ failles orientées et de $n$ homéomorphismes de recollement, le cas $n = 2$
étant écrit à la page 8.

\pagerange{10}{13}
\section*{Existence, et conditions sur les lieux (pages 10 à 13)}

\subsection*{Quels espaces portent une faille}

\begin{quote}
\emph{Un espace compact triangulé admet une structure de faille réduite si et
seulement s'il est connexe et non réduit à un point.}
\end{quote}

« C'est gros ! », dit la page, et c'est juste\footnote{La page ne démontre rien.
Pour $X$ connexe, les deux points $a$, $b$ doivent être des points de
non-coupure ; tout continu non réduit à un point en a au moins deux (R.~L.~Moore,
1920), et dans un espace triangulé tout point a un voisinage conique, le cône sur
son entrelacs.}. La page ajoute qu'il y en a alors une infinité, et que la
structure est unique pour le segment, et pour lui seul. La seconde affirmation est
exactement le théorème de caractérisation de l'arc : un continu métrique qui a
exactement deux points de non-coupure est un arc. La première demande une
réserve : un arbre fini à $k$ feuilles n'a que $k(k-1)/2$ structures de faille
réduite, une par paire de feuilles\footnote{La page écrit « $\infty$ de
structures de faille réduite ». C'est vrai dès que $X$ contient un cycle ou une
cellule de dimension $\geq 2$ ; faux pour les arbres. Les structures \emph{non}
réduites, elles, sont toujours en infinité : la marge donne sur $[0,1]$ la famille
$A = [0,\alpha]$, $B = [\beta, 1]$, $0 \leq \alpha < \beta \leq 1$, qui ne
contredit pas l'unicité de la structure réduite.}.

\subsection*{Deux essais retirés}

La page considère la relation d'équivalence entre structures de faille orientée
sur $X$ engendrée par : $(A,B) \sim (A',B')$ si $A \subset A'$ et $B \subset B'$.
Elle affirme d'abord que si $(X, \{A,B\})$ est une faille et $A \subset A'$,
$B \subset B'$ sont fermés, régulièrement immergés et disjoints, alors
$(X, \{A',B'\})$ est une faille --- et barre l'énoncé, avec en marge « Faux !
idiot ! » : grossir un bout peut déconnecter le complémentaire de l'autre. Elle
affirme ensuite que deux structures de faille orientée sur un $X$ triangulé sont
toujours équivalentes, et note en marge « Faux (cas orienté !) déjà dans le cas des
segments ! ». En effet, sur $[0,1]$, chaque pas de la relation garde $0$ dans
l'origine ; les orientations $([0,\alpha], [\beta,1])$ et $([\beta,1],
[0,\alpha])$ ne sont pas équivalentes.

\subsection*{Imposer une propriété aux bouts et aux lieux}

On peut exiger de $A$, $B$ et des lieux une propriété $P$ intrinsèque --- être
connexes, contractiles, homéomorphes à un espace donné --- et ne considérer que
l'ensemble $\mathcal{L}_{P}$ des lieux qui la possèdent. Pour que $\mathcal{L}_{P}$
garde les propriétés a) à c) de la page 9, il est bon d'imposer $P$ aussi aux bords
$S_{A}$, $S_{B}$ des voisinages coniques des bouts, et, pour un lieu intérieur $C$,
aux bords de ses voisinages coniques dans $X_{A,C}$ et dans $X_{B,C}$. Exemples :
\begin{enumerate}
\item[1)] \emph{Être réduit à un point.} On retrouve la théorie des segments
topologiques ordonnés : les lieux sont les points intérieurs, l'ordre est celui du
segment.
\item[2)] \emph{Être homéomorphe à $B^{n-1}$}\footnote{La page annonce « être
homéomorphe à $B^{n}$ », mais ce qui est homéomorphe à $B^{n}$ est l'espace $X$ ;
les bouts et les lieux du croquis, dessiné pour $n = 2$, sont des $B^{1}$ (le
croquis écrit $D^{1}$ pour l'un d'eux).}. On trouve la théorie de la décomposition
de $I \times B^{n-1} \simeq B^{n}$, de bouts $\{0\} \times B^{n-1}$ et $\{1\} \times
B^{n-1}$, par des lieux homéomorphes à $B^{n-1}$ ; ou, en contractant les bouts,
d'une boule munie de deux points de son bord.
\item[3)] \emph{Être homéomorphe à $S^{n}$.} Sur $S^{n+1}$ munie de deux points
distincts $a$, $b$, on trouve la théorie de la décomposition par des sphères
$S^{n}$ ; pour $n = 0$, un cercle coupé par des paires de points.
\end{enumerate}

On peut enfin exiger qu'un lieu intérieur $C$ ait un voisinage homéomorphe à
$C \times {]-1,1[}$, avec $C = C \times \{0\}$, et les deux « rives »
$C \times {]-1,0[}$ et $C \times {]0,1[}$ contenues respectivement dans $X_{A,C}$ et
$X_{B,C}$. C'est la notion de \emph{plongement bicollier}\footnote{M.~Brown (1962)
a montré qu'un sous-espace localement bicollier est bicollier ; son théorème de
Schoenflies généralisé (1960, et B.~Mazur 1959) dit qu'une sphère $S^{n}$
bicollier dans $S^{n+1}$ borde deux boules. C'est exactement ce qui fait marcher
l'exemple 3). Le mot « rives » apparaît ici pour la première fois ; il devient
une notion à la page 14.}.

\pagerange{14}{15}
\section*{Rives et bi-rives (pages 14 et 15)}

Cette partie ouvre, le 7 juin, un deuxième paragraphe intitulé « Boucle
topologique »\footnote{Le titre est souligné de deux traits, dont l'un pourrait
être une rature ; la page enchaîne aussitôt sur un autre titre, « Rives d'une
partie fermée ».}.

Pour un espace $Z$, soit
\[
\mathrm{Comp}(Z) = \mathrm{Hom}_{\mathrm{cont}}(Z, \{0,1\}) = \Gamma(Z,
\mathbf{F}_{2,Z}),
\]
l'anneau booléen des parties ouvertes et fermées de $Z$. Pour une partie fermée $A$
de $X$, d'ouvert complémentaire $j : U = X \smallsetminus A \hookrightarrow X$, on
pose
\[
\mathrm{Riv}(A,X) = \varinjlim_{V \supset A} \mathrm{Comp}(V \smallsetminus A) =
\varinjlim_{V \supset A} \Gamma(V, j_{*}\mathbf{F}_{2,U}),
\]
la limite inductive portant sur les voisinages $V$ de $A$\footnote{La page note
$i$ l'inclusion de $U$ ; on écrit $j$ pour ne pas la confondre avec une
immersion fermée.}. Ses éléments sont les \emph{rives} de $A$ : des germes, le
long de $A$, de parties ouvertes et fermées de $X \smallsetminus A$.

\emph{Si $A$ est régulièrement immergé}, les $W_{t} \smallsetminus A \simeq
S_{A} \times {]0,t]}$ sont cofinaux, et
\[
\mathrm{Riv}(A,X) \simeq \mathrm{Comp}(S_{A}).
\]

\emph{Si $X$ est paracompact}, pour tout faisceau $\mathcal{G}$ sur $X$ et tout
fermé $A$, $\Gamma(A, \mathcal{G}|_{A}) = \varinjlim_{V \supset A}
\mathcal{G}(V)$\footnote{C'est le théorème de R.~Godement (\emph{Topologie
algébrique et théorie des faisceaux}, 1958, II.3.3.1), que la page cite sans
référence.}.

\emph{Le découpage.} Soit $\widetilde{A} \to A$ le spectre de Stone relatif du
faisceau d'anneaux booléens $(j_{*}\mathbf{F}_{2,U})|_{A}$ : au-dessus de $x \in A$,
ses points sont ceux du spectre de la fibre, c'est-à-dire, grossièrement, les côtés
de $X \smallsetminus A$ au voisinage de $x$\footnote{La page écrit $\widetilde{A} =
\mathrm{Spec}(i_{*}\mathbf{F}_{2,U})$ sans préciser ; la lecture comme spectre de
Stone relatif est la nôtre, et c'est celle qui rend vraie l'égalité suivante, par
la dualité de Stone (1936) fibre à fibre.}. On a
\[
\mathrm{Comp}(\widetilde{A}) \simeq \Gamma(A, (j_{*}\mathbf{F}_{2,U})|_{A}),
\]
d'où une application canonique
\[
(*) \qquad \mathrm{Riv}(A,X) \longrightarrow \mathrm{Comp}(\widetilde{A}),
\]
toujours injective, et bijective si $X$ est paracompact\footnote{L'injectivité
n'utilise rien : une section sur $V$ nulle en germe en chaque point de $A$ est
nulle sur un voisinage ouvert de $A$. La bijectivité est le théorème de Godement
ci-dessus. La page écrit d'abord « bijective », puis ajoute « injective, et » au
dessus de la ligne.}. Dans ce cas, choisir une rive, c'est choisir une partie
ouverte et fermée de $\widetilde{A}$. On pense à $\widetilde{A}$ comme au bord de
l'espace obtenu en coupant $X$ le long de $A$ : pour une hypersurface bilatère,
$\widetilde{A} = A \times \{\pm 1\}$ ; pour l'âme d'un ruban de Möbius,
$\widetilde{A}$ est le revêtement connexe de degré 2 du cercle.

\begin{quote}
\textbf{Définition} (page 15). Une \emph{bi-rive} de $A$ est une paire
$\beta = \{\rho, \rho'\}$ d'éléments de $\mathrm{Riv}(A,X)$ tels que
$\rho + \rho' = 1$ --- une décomposition en deux morceaux de $X \smallsetminus A$ au
voisinage de $A$. Elle est \emph{triviale} si $\{\rho, \rho'\} = \{0,1\}$. On dit que
$A$ \emph{déconnecte localement} $X$ s'il existe une bi-rive non triviale.
\end{quote}

Dans un anneau booléen, $\rho = \rho'$ et $\rho + \rho' = 1$ imposent $1 = 0$, ce
qui arrive seulement si $A$ est ouvert ; une bi-rive a donc deux éléments distincts
dès que $A$ n'est pas ouvert\footnote{La page écrit « si $A \neq \emptyset$ ». Pour
$A$ ouvert et fermé non vide, $V = A$ est un voisinage et $\mathrm{Riv}(A,X) = 0$.}.
Une bi-rive découpe $\widetilde{A} = \widetilde{A}_{1} \sqcup \widetilde{A}_{2}$.
Pour une hypersurface bilatère connexe, $\mathrm{Riv}(A,X)$ a quatre éléments, $0$,
$1$ et les deux côtés, qui forment une bi-rive non triviale. L'âme $A$ du ruban de
Möbius coupe le ruban en deux près de chacun de ses points, mais le bord de son
voisinage tubulaire est un cercle connexe : $\mathrm{Riv}(A,X) = \{0,1\}$, et il n'y
a pas de bi-rive non triviale. C'est une hypersurface unilatère.

\pagerange{16}{18}
\section*{Mailles (pages 16 à 18)}

Soit $A$ fermé dans $X$ et $\beta = \{\rho, \rho'\}$ une bi-rive. On note
$\widetilde{X}$ l'espace $X$ coupé le long de $A$, qui contient $\widetilde{A}$ et
s'envoie sur $X$ en envoyant $\widetilde{A}$ sur $A$\footnote{La page 16 se sert de
$\widetilde{X}$ sans le définir ; la page 14 écrit seulement « la paire
$(\widetilde{X}, \widetilde{A})$ ». La définition précise, en recollant $X
\smallsetminus A$ et $\widetilde{A}$ par la topologie des rives, n'est pas sur la
page et n'est pas faite ici.}. On contracte ensuite chacune des deux moitiés
$\widetilde{A}_{1}$, $\widetilde{A}_{2}$ sur $A$ :
\[
\overline{X} = \mathrm{Dec}_{\beta}(A, X) = \widetilde{X}
\amalg_{\widetilde{A}_{1}} A \amalg_{\widetilde{A}_{2}} A .
\]
C'est le \emph{découpage} de $X$ par $\beta$. Il contient une partie canoniquement
isomorphe à $A \times \beta$, réunion de deux copies $A_{\rho}$, $A_{\rho'}$ de $A$,
et l'application $\overline{X} \to X$ induit un homéomorphisme
\[
\overline{X} \smallsetminus (A_{\rho} \cup A_{\rho'}) \xrightarrow{\ \simeq\ } X
\smallsetminus A
\]
et, sur $A \times \beta$, la projection $A \times \beta \to A$.

\begin{quote}
\textbf{Définition} (pages 16 et 17). On dit que $(A, \beta)$ définit sur $X$ une
structure de \emph{maille} de \emph{nœud} $A$ et de bi-rive $\beta$ si
$(\overline{X}, \{A_{\rho}, A_{\rho'}\})$ est une faille topologique. Orienter la
maille, c'est orienter cette faille, c'est-à-dire choisir $\rho$ ou $\rho'$.
\end{quote}

La page 16 formulait d'abord la condition sur $(\widetilde{X}, \{\widetilde{A}_{1},
\widetilde{A}_{2}\})$, avec « $A$ régulièrement immergé dans $X$ », en notant que
$\widetilde{A}$ l'est alors dans $\widetilde{X}$\footnote{Un premier nom, « bouclage
régulier », est biffé au profit de « maille ». La marge dit « réciproque » devant la
remarque qu'une paire satisfaisant a), b), c) est une faille dès que
$\widetilde{A}$ est régulièrement immergé.}. Une maille est ainsi un espace obtenu
en recollant les deux bouts d'une faille par un homéomorphisme : le cercle, obtenu
d'un segment ; l'anneau, obtenu d'un carré en recollant deux côtés opposés, de
nœud un rayon. En termes actuels, le nœud d'une maille est une hypersurface
bilatère non séparante\footnote{Le rapprochement est le nôtre. Pour une variété
$X$, une telle hypersurface est la fibre en une valeur régulière d'une application lisse
$X \to S^{1}$, et définit une classe de $H^{1}(X; \mathbf{Z})$ ; la notion de
maille ne demande aucune de ces structures.}.

\subsection*{Deux mailles compatibles}

Soient $(A_{1}, \beta_{1})$ et $(A_{2}, \beta_{2})$ deux structures de maille sur
$X$, $\beta_{i} = \{\rho_{i}, \rho'_{i}\}$. On dit qu'elles sont \emph{compatibles}
si :
\begin{enumerate}
\item[a)] $A_{1} \cap A_{2} = \emptyset$, de sorte que $A_{2}$ s'identifie à une
partie de $\overline{X}^{(A_{1}, \beta_{1})} \smallsetminus A_{1} \times
\beta_{1}$ ;
\item[b)] $A_{2}$ est un lieu de la faille $\overline{X}^{(A_{1}, \beta_{1})}$ ;
\item[c)] si $U$ et $V$ sont les deux parties de $\overline{X}^{(A_{1},\beta_{1})}
\smallsetminus (A_{1} \times \beta_{1} \cup A_{2})$ que découpe ce lieu, la paire
$\{U, V\}$ définit à la fois la bi-rive $\beta_{1}$ de $A_{1}$ et la bi-rive
$\beta_{2}$ de $A_{2}$\footnote{La page parle des « deux composantes connexes » ; par
définition d'un lieu, ce sont les deux termes de la somme topologique, qui ne sont
pas nécessairement connexes. La condition c) commence au bas de la page 17 et
s'achève en haut de la page 18.}.
\end{enumerate}

Le croquis de la page 17 est un anneau coupé par deux rayons, l'un en haut de nœud
$A$, l'autre en bas, les deux moitiés étant $U$ et $V$. La page 18 affirme que cette
condition est en fait symétrique en les deux nœuds $(A_{i}, \beta_{i})$, et le
texte s'arrête là\footnote{La symétrie n'est pas démontrée. Elle est plausible :
$U$ et $V$, vus dans $X \smallsetminus (A_{1} \cup A_{2})$, ne dépendent pas de
l'ordre dans lequel on coupe ; mais b) pour $A_{1}$ dans
$\overline{X}^{(A_{2},\beta_{2})}$ reste à vérifier.}.

\end{document}
